\section{Related Work}

\subsection{Security Vulnerabilities in LLM Agents}
The integration of LLMs with external tools introduces a new attack surface, primarily driven by Prompt Injection (PI) attacks. Greshake et al. demonstrated "Indirect Prompt Injection," where an LLM processing untrusted content (e.g., a website) can be manipulated to execute malicious instructions \cite{greshake2023}. Liu et al. further showed that standard safety alignment (RLHF) is often insufficient to prevent "jailbreaking" when models are placed in an agentic loop \cite{liu2023}.

Unlike standalone chatbots, agents have the capability to execute code and modify file systems. The OWASP Top 10 for LLM Applications highlights "Insecure Plugin Design" and "Excessive Agency" as critical risks \cite{owasp2023}. A compromised agent effectively becomes a "confused deputy," leveraging the user's credentials to perform unauthorized actions.

\subsection{Existing Defense Mechanisms}
Current defenses largely fall into two categories: static guardrails and dynamic monitoring.
Static approaches, such as NeMo Guardrails \cite{nemo}, use predefined rules to filter inputs and outputs. While effective against known patterns, they lack the semantic understanding to detect context-dependent attacks (e.g., distinguishing between a user asking to delete a temporary file versus a system file).

Dynamic approaches leverage "LLM-as-a-Judge" to evaluate inputs. However, early implementations often suffered from high latency and cost. Recent works have proposed lightweight "self-correction" mechanisms \cite{pan2023}, but these often rely on the same model for both execution and verification, leaving them vulnerable to single-point failures if the model is successfully jailbroken.

\subsection{Zero Trust Architecture (ZTA) for AI}
Zero Trust, as defined by NIST SP 800-207 \cite{nist}, posits that trust should never be implicit. While ZTA has been widely adopted in network security (e.g., Google's BeyondCorp \cite{beyondcorp}), its application to AI agents is nascent.
Recent guidelines from BSI and ANSSI \cite{bsi_anssi} argue for applying ZTA principles---specifically Least Privilege and Continuous Verification---to the AI model itself.

Our work, \texttt{llm-cli}, bridges these domains. Unlike existing agent frameworks that focus on capability (e.g., AutoGen \cite{autogen}), we prioritize security by implementing a distinct workload identity and an independent, secondary LLM for intent verification, effectively decoupling the "actor" from the "auditor."
